\documentclass[11pt]{article}
\usepackage[margin=1in]{geometry}
\usepackage{amsmath,amssymb,amsthm}
\usepackage{hyperref}

\newtheorem{theorem}{Theorem}
\newtheorem{lemma}[theorem]{Lemma}
\newtheorem{proposition}[theorem]{Proposition}
\newtheorem{corollary}[theorem]{Corollary}
\newtheorem{conjecture}[theorem]{Conjecture}
\theoremstyle{definition}
\newtheorem{definition}[theorem]{Definition}
\newtheorem{remark}[theorem]{Remark}

\newcommand{\Z}{\mathbb{Z}}
\newcommand{\R}{\mathbb{R}}
\newcommand{\eps}{\varepsilon}

\title{A magnitude lower bound and a lattice-clustering reformulation \\
of the Umans--Wang $(\alpha,\beta)$-divisor conjecture}
\author{Eric Schultz\\ \small independent researcher \\ \small ORCID 0009-0006-6283-1696}
\date{July 2026}

\begin{document}
\maketitle

\begin{abstract}
The Umans--Wang $(\alpha,\beta)$-divisor conjecture (arXiv:2511.10851) asks, for
fixed $\alpha,\beta<1/2$, whether there exist sets $S,T\subseteq\Z^+$ with
$|S|,|T|\le n^{\beta}$ and elements at most $\exp(n^{\alpha})$ such that every
$i\in[n]$ divides $s-t$ for some $s\in S,\,t\in T$; a positive answer would break
the exponent-$3/2$ barrier for polynomial factorization. We record two things.
First, an elementary but self-contained \emph{magnitude lower bound}
$2\beta+\alpha\ge 1$, obtained via the discriminant of the covering set; this is
consistent with the conjecture and reproduces, by an independent route, the bound
implied by the primorial. Second, and our main contribution, we reduce the
conjecture at fixed $\alpha$ to a precise \emph{lattice-clustering} statement: the
covering set is a point of a congruence lattice of covolume $\approx e^{n/2}$ that
must lie inside an axis-parallel box of volume $e^{K n^{\alpha}}$, and a
generic lattice of that covolume meets the box if and only if $\alpha+\beta\ge 1$.
Since the counting floor sits at $\alpha+\beta=(1+\alpha)/2<1$, a solution requires
a \emph{special} lattice whose points cluster in an exponentially small box. We show
this clustering property is \emph{equivalent} to the conjecture (not weaker), and
observe that clustering of this \emph{type} is what the bounded-root-discriminant
unit lattices of the 2026 refutation of the real sum--product conjecture by
Bloom--Sawin--Schildkraut--Zhelezov provide. We are careful about the degree
scaling: our setting appears to require such fields at \emph{polynomial} degree
$d\asymp n^{(1-\alpha)/2}$, whereas their construction operates at
\emph{logarithmic} degree, so their result does not directly supply what we need;
the requisite fields exist at every degree, but whether the shortness carries over
at polynomial degree is part of the open question, not something the reduction
settles. This isolates a single, frontier-adjacent open question and rules out a
family of would-be elementary obstructions.
\end{abstract}

\section{Introduction}

Umans and Wang \cite{UW} reduce beating the longstanding exponent-$3/2$ barrier for
distinct-degree factorization of degree-$n$ polynomials over $\mathbb{F}_q$ to a
number-theoretic conjecture about covering $[n]$ by divisors of a difference set.

\begin{conjecture}[$(\alpha,\beta)$-divisor conjecture \cite{UW}]\label{conj:uw}
For infinitely many $n$ there exist $S,T\subseteq\Z^{+}$ such that
$A=\{s-t : s\in S,\ t\in T\}$ satisfies the $n$-divisor property
(every $i\in\{1,\dots,n\}$ divides some element of $A$), with
(i) all elements of $S,T$ of magnitude at most $\exp(n^{\alpha})$, and
(ii) $|S|,|T|\le n^{\beta}$.
\end{conjecture}

Achieving $\alpha+\beta\le 1+o(1)$ is trivial; the content is $\alpha,\beta<1/2$,
and $\alpha=\beta=1/3$ would yield an exponent-$4/3$ factoring algorithm \cite{UW}.
Throughout we work at a \emph{fixed} $\alpha<1/2$ and ask for the smallest feasible
$\beta$.

\paragraph{Reduction to the primes in $(n/2,n]$.}
The $n$-divisor property is governed by its hardest instances, the primes
$p\in P(n):=\{p \text{ prime}: n/2<p\le n\}$: any $i\le n$ has all prime-power
factors at most $n$, and the primes in $(n/2,n]$ are those that can divide at most
one difference of magnitude $O(n)$ and cannot share factors. We therefore study
covering $P(n)$ by divisibility. By the prime number theorem
$\log\prod_{p\in P(n)}p=\theta(n)-\theta(n/2)=\tfrac{n}{2}(1+o(1))$
(verified numerically: the ratio is $0.963,0.997,0.999$ at $n=10^3,10^4,10^5$).

\section{The magnitude lower bound}

We first give a clean lower bound. Write $M=\exp(n^{\alpha})$ and, for a covering
set $S$ with $|S|=K=n^{\beta}$, let
$D(S)=\prod_{i<j}(s_i-s_j)$ be the Vandermonde product (the square root of the
discriminant of $\prod_i(x-s_i)$).

\begin{lemma}\label{lem:divides}
If the difference set of $S$ covers $P(n)$ by divisibility, then
$Q:=\prod_{p\in P(n)}p$ divides $D(S)$.
\end{lemma}
\begin{proof}
Each $p\in P(n)$ divides some difference $s_i-s_j$, hence divides $D(S)$. The
primes of $P(n)$ are distinct, so their product $Q$ divides $D(S)$.
\end{proof}

We verified Lemma~\ref{lem:divides} directly on explicit covers at $n=20,30,40$
(e.g.\ $S=\{0,17,20,39\}$ covers $P(20)$ and $Q=46189\mid D(S)$).

\begin{theorem}\label{thm:lb}
Any $S,T$ satisfying Conjecture~\ref{conj:uw} obey $2\beta+\alpha\ge 1-o(1)$.
\end{theorem}
\begin{proof}
Take $S$ with $|S|=K=n^{\beta}$ (the symmetric case; the asymmetric $S-T$ case is
identical with $K^2$ in place of $\binom{K}{2}$). By Lemma~\ref{lem:divides},
$Q\mid D(S)$, so $\log Q\le \log|D(S)|$. Each of the $\binom{K}{2}$ factors of
$D(S)$ has magnitude at most $2M$, hence
\[
\tfrac{n}{2}(1-o(1))=\log Q\le \binom{K}{2}\log(2M)\le \tfrac{K^2}{2}\,n^{\alpha}.
\]
Thus $K^2 n^{\alpha}\ge n(1-o(1))$, i.e.\ $2\beta+\alpha\ge 1-o(1)$.
\end{proof}

\begin{remark}
Theorem~\ref{thm:lb} is consistent with Conjecture~\ref{conj:uw}: at
$\alpha=\beta=0.4$ one has $2\beta+\alpha=1.2\ge 1$. It reproduces, via the
discriminant, the bound that Umans--Wang obtain from
$\prod_{p<n}p=e^{n(1+o(1))}$. It is tight \emph{as a magnitude bound}: it cannot by
itself yield a stronger exponent, and in particular cannot decide the conjecture.
The boundary $2\beta+\alpha=1$ is the \emph{counting floor}
$\beta=(1-\alpha)/2$.
\end{remark}

\section{The lattice-clustering reformulation}

We now reduce the conjecture at fixed $\alpha$ to a single geometric statement. Fix
$\alpha<1/2$ and set $K=n^{(1-\alpha)/2}$ (the counting floor), $B=\exp(n^{\alpha})$.

\paragraph{The congruence lattice.}
A covering assignment attaches to each prime $p\in P(n)$ a pair of indices whose
difference $p$ must divide. Collecting these, the admissible tuples
$s=(s_1,\dots,s_K)$ form a coset structure inside the lattice
\[
L=\{s\in\Z^{K}: s_i\equiv s_j \pmod{m_{ij}}\ \text{for each assigned pair}\},
\]
where $m_{ij}$ is the product of primes assigned to pair $(i,j)$. Since each prime
of $P(n)$ is assigned to exactly one pair, the moduli are pairwise coprime and the
index $[\Z^K:L]=\prod_{p\in P(n)}p=Q$, so $L$ has covolume $Q=e^{n/2(1+o(1))}$ in
$\R^K$.

\paragraph{The box condition.}
A covering set of magnitude $\le B$ with distinct coordinates is exactly a point of
$L$ in the box $[0,B]^K$ with distinct coordinates. The box has volume
$B^{K}=\exp(K n^{\alpha})$. Comparing to the covolume:
\[
\frac{\operatorname{vol}[0,B]^K}{\operatorname{covol}(L)}
=\exp\!\big(K n^{\alpha}-\tfrac{n}{2}\big),
\qquad
K n^{\alpha}=n^{(1-\alpha)/2}\cdot n^{\alpha}=n^{(1+\alpha)/2}.
\]

\begin{proposition}\label{prop:generic}
For a generic lattice of covolume $Q$, the expected number of points in $[0,B]^K$
is $B^K/Q=\exp\!\big(n^{(1+\alpha)/2}-\tfrac n2\big)$. This is
$\ge 1$ if and only if $K n^{\alpha}\ge \tfrac n2$, i.e.\ (writing $K=n^{\beta}$)
if and only if $\alpha+\beta\ge 1$.
\end{proposition}

At the counting floor $\beta=(1-\alpha)/2$ we have
$\alpha+\beta=(1+\alpha)/2<1$, so the box is exponentially smaller than the
covolume and a \emph{generic} congruence lattice contains \emph{no} admissible
point (numerically, $B^K/Q=e^{-\Theta(n)}$; the ratio $Kn^{\alpha}/\log Q$ is
$0.017$ and $0.032$ at $n=10^6$, $\alpha=0.3,0.4$). Hence:

\begin{theorem}[Reformulation]\label{thm:reform}
Conjecture~\ref{conj:uw} at fixed $\alpha$ is equivalent to the existence of
congruence lattices $L$ as above whose points \emph{cluster}: $L$ must contain a
distinct-coordinate point in a box that is a factor $e^{\Theta(n)}$ smaller in
volume than $\operatorname{covol}(L)$. Equivalently, $L$ must beat the generic
point-count by an $e^{\Theta(n)}$ factor.
\end{theorem}

\begin{proof}[Proof sketch]
A solution to Conjecture~\ref{conj:uw} at $\alpha$ (with $\beta$ at the floor) is,
after fixing an assignment of primes to pairs, precisely a distinct-coordinate
point of the corresponding $L$ in $[0,B]^K$; conversely such a point yields a
covering set of the required magnitude and cardinality. The magnitude/cardinality
bookkeeping is Theorem~\ref{thm:lb} and Proposition~\ref{prop:generic}. The freedom
to choose the assignment does not change the covolume $Q$, only the shape of $L$;
so feasibility is exactly the existence of \emph{some} assignment whose lattice
clusters into the box.
\end{proof}

\begin{remark}[The reduction does not simplify the problem, and that is the point]
Theorem~\ref{thm:reform} is an \emph{equivalence}, not a weakening. Any attempt to
prove the clustering as a lighter sub-lemma is therefore misguided: it carries the
full weight of the conjecture. The value of the reformulation is the change of
language---from divisor covering to lattice clustering---which connects the problem
to a concrete construction technology (below) and, as we discuss in the companion
note, explains why an entire family of elementary ``barrier'' arguments fail: they
attempt to lower-bound coordinates of $L$ by a generic (Minkowski-average) estimate,
which the clustering hypothesis is precisely designed to beat.
\end{remark}

\section{The construction frontier}

The clustering property in Theorem~\ref{thm:reform} is not abstract: it is exactly
the property exploited by Bloom, Sawin, Schildkraut and Zhelezov \cite{BSSZ} in
their 2026 refutation of the real sum--product conjecture. They construct sets
$A=GP$ in totally real number fields $K$ of degree $d\asymp\log|A|$ with
bounded root discriminant (Martinet towers), where $G$ is a box in the unit lattice
(multiplicatively structured, small product set) and $P$ a box of algebraic
integers (additively structured). In such a field the unit lattice has
$O(1)$-size-per-dimension vectors despite large covolume---exactly the
super-generic clustering our reformulation requires.

We record the quantitative situation, and flag honestly where it is favorable and
where it is not. To place $K=n^{(1-\alpha)/2}$ coordinates below $B=e^{n^{\alpha}}$
against covolume $Q=e^{n/2}$, one needs shortness beating the Minkowski average
$Q^{1/K}=e^{n^{(1+\alpha)/2}/2}$ by a factor growing polynomially in $n$
(numerically $9\times,58\times,316\times$ at $\alpha=0.3$, $n=10^4,10^6,10^8$).
A degree-$d$ bounded-root-discriminant field has unit lattice of rank $d-1$ with
$O(1)$-size-per-dimension vectors, which is shortness of the required \emph{type};
and Martinet towers provide bounded-root-discriminant fields at \emph{every} degree
$d$, so in particular the fields exist at any degree we might ask for.

\emph{However, the degree scaling is a genuine soft point of this reduction, and we
do not claim it is resolved.} Placing $K$ coordinates subject to $K$ constraints
calls, naively, for on the order of $K$ independent short directions, hence degree
$d\asymp K=n^{(1-\alpha)/2}$---\emph{polynomial} in $n$. This is a different regime
from the construction of \cite{BSSZ}, whose degree scales \emph{logarithmically},
$d\asymp\log|A|$, in the size of the set built; the ratio $K/\log n$ grows without
bound (it exceeds $500$ already at $\alpha=0.3$, $n=10^{12}$). Thus one cannot
simply cite \cite{BSSZ} for the shortness we need: their $O(1)$-per-dimension bound
is established for sets of size $|A|$ at degree $\log|A|$, whereas using degree
$d\asymp K$ corresponds to an object of size $|A|\asymp e^{K}$, larger than and
distinct from what they analyze. What the Martinet input gives us is the
\emph{existence} of the fields at the required (polynomial) degree; whether their
unit lattices supply $K$ \emph{usable, independent} short directions respecting our
specific pair-congruences is not delivered by \cite{BSSZ} and is, in fact, part of
the open question below.

Consequently the open question is twofold and precise: whether the specific
pair-congruences of $L$ can be embedded into such a unit lattice
(i)~\emph{at the required polynomial degree} $d\asymp K$, and (ii)~\emph{preserving
$O(1)$-per-dimension shortness}, so that a simultaneous solution descends to $K$
distinct integers $\le B$. This is a \emph{descent-to-$\Z$} question of the type
\cite{BSSZ} leave open (they resolve the $\R$, $p$-adic, and finite-field cases;
the integer case, and hence the sum--product problem in $\Z$, remains open), now
with the additional, unresolved demand that the shortness survive at polynomial
rather than logarithmic degree.

\paragraph{Structural notes for a specialist.}
Three facts constrain any attack.
(1) The conjecture uses an \emph{asymmetric} difference set $S-T$; the two-set
freedom is genuine (exact search shows it beats the naive per-slot magnitude bound
by a growing factor, verified at $K\le 3$), and it matches the $G\cdot P$ split of
\cite{BSSZ}, with $S$ playing the role of the multiplicative box and $T$ the
additive box. (2) The relevant regime is \emph{bundled}: at the floor each
difference must be divisible by $\Theta(n^{\alpha}/\log n)$ primes of $P(n)$
simultaneously, forcing magnitude $\approx e^{n^{\alpha}}$; small-magnitude or
polynomial-magnitude computation cannot probe it (a difference $\le 2n$ carries at
most one prime $>n/2$). (3) A solution must be \emph{structured} (efficiently
computable) for the algorithmic application, so an existence proof alone does not
suffice; the explicit \cite{BSSZ}-style construction is the right kind of target.

\section{Conclusion}

We have isolated the Umans--Wang conjecture at fixed $\alpha$ as a lattice-clustering
problem equivalent to the conjecture itself, and tied it to the concrete
bounded-root-discriminant construction technology of \cite{BSSZ}. The magnitude
bookkeeping is favorable and the requisite fields exist at every degree; but the
degree scaling is not resolved---our setting appears to demand shortness at
\emph{polynomial} degree $d\asymp n^{(1-\alpha)/2}$, while \cite{BSSZ} supply it at
\emph{logarithmic} degree---so the reduction leaves a descent-to-$\Z$ embedding
question that is genuinely open and includes this degree gap, rather than a merely
technical residue. The proven bound $2\beta+\alpha\ge 1$ is the only unconditional
statement; the conjecture remains open in both directions. What we can claim is
narrower but solid: the elementary obstructions to the conjecture provably fail
(they are artifacts of generic estimates), and the natural constructive route is now
explicit and precisely delimited, including an honest account of where it is not yet
supported.

\begin{thebibliography}{9}
\bibitem{UW} C.~Umans and S.~Wang. \emph{A number-theoretic conjecture implying
faster algorithms for polynomial factorization and integer factorization.}
arXiv:2511.10851, 2025.
\bibitem{BSSZ} T.~F.~Bloom, W.~Sawin, C.~Schildkraut, and D.~Zhelezov.
\emph{The sum-product conjecture is false for real numbers.} arXiv:2605.28781, 2026.
\bibitem{ABP} R.~Agrawal, T.~F.~Bloom, and G.~Petridis. \emph{More on the
sum-product problem for integers with few prime factors.} arXiv:2512.04931, 2025.
\end{thebibliography}

\end{document}
